\subsection{Final Experimental Design}

The last experiment was executed based on four practical scenarios that were aimed at testing both outdoor and indoor conditions that would apply to the last prototype in real use. The tests in the outdoor environment were carried out at the distances of 20~m and 30~m,  whereas indoor ones were performed at 10 m distance and with and without the line of sight between two devices. All three alignment strategies were tested under equal search range, RSSI collection conditions, and mechanical accuracy.

All the practical scenarios were run ten times for each alignment strategy. The data for RSSI, convergence time, number of movements, and success rates mentioned in the following subsections were calculated based on the average results from those runs.

\setlength{\tabcolsep}{4pt}

\vspace{0.5cm}
\begin{table}[H]
\centering
\caption{Shared Runtime Parameters for Practical Trials}
\label{tab:practical_runtime_results}
\renewcommand{\arraystretch}{1.25}
\begin{adjustbox}{max width=\textwidth}
\begin{tabular}{|
>{\raggedright\arraybackslash}p{1.7cm}|
>{\raggedright\arraybackslash}p{1.7cm}|
>{\raggedright\arraybackslash}p{1.4cm}|
>{\raggedright\arraybackslash}p{1.4cm}|
>{\raggedright\arraybackslash}p{1.4cm}|
>{\raggedright\arraybackslash}p{1.8cm}|
>{\raggedright\arraybackslash}p{2.0cm}|
>{\raggedright\arraybackslash}p{1.8cm}|}
\hline
\textbf{Pan Range} & \textbf{Tilt Range} & \textbf{Pan Step} & \textbf{Tilt Step} & \textbf{Fine Step} & \textbf{Steps/ Deg.} & \textbf{Samples/ Point} & \textbf{Sample Delay} \\
\hline
$0^\circ$--$180^\circ$ & $25^\circ$--$75^\circ$ & $10^\circ$ & $5^\circ$ & $2^\circ$ & 5 & 20 & 10~ms \\
\hline
\end{tabular}
\end{adjustbox}
\end{table}

\vspace{0.5cm}
\begin{table}[H]
\centering
\caption{Practical Test Scenarios Used in Final Evaluation}
\label{tab:practical_scenarios}
\renewcommand{\arraystretch}{1.25}
\begin{adjustbox}{width=\textwidth}
\begin{tabular}{|
>{\raggedright\arraybackslash}p{1.15cm}|
>{\raggedright\arraybackslash}p{1.85cm}|
>{\raggedright\arraybackslash}p{1.35cm}|
>{\raggedright\arraybackslash}p{3.85cm}|
>{\raggedright\arraybackslash}p{4.25cm}|}
\hline
\textbf{ID} & \textbf{Env.} & \textbf{Dist.} & \textbf{Channel Condition} & \textbf{Purpose} \\
\hline
S1 & Outdoor & 20~m & Clear LOS & Reference case with minimal obstruction and strong received power \\
\hline
S2 & Outdoor & 30~m & Partial foliage & Tests longer range and mild attenuation in a mostly open channel \\
\hline
S3 & Indoor & 10~m & Corridor LOS with reflections & Evaluates sensitivity to multipath in a confined environment \\
\hline
S4 & Indoor & 10~m & Wall and corridor corner & Hardest indoor case with obstruction, fading, and multipath peaks \\
\hline
\end{tabular}
\end{adjustbox}
\end{table}

\subsection{Practical Trial Results}

The practical trial data serves as the main evidence source for the concluding experiment, since practical experiments always contain such parameters as RF noise, actuator delay and mechanical inaccuracies that cannot be completely reflected in the simulation experiment. Here, in the presented table, the reported RSSI standard deviation refers to the RSSI variability index from the methodology and is used in this sense for representing repeatability of the process over several trial runs.

\vspace{0.5cm}
\begin{table}[H]
\centering
\caption{Average Practical Performance Across All Scenarios}
\label{tab:overall_practical_performance}
\renewcommand{\arraystretch}{1.25}
\begin{adjustbox}{width=\textwidth}
\begin{tabular}{|
>{\raggedright\arraybackslash}p{2.35cm}|
>{\raggedright\arraybackslash}p{2.0cm}|
>{\raggedright\arraybackslash}p{1.9cm}|
>{\raggedright\arraybackslash}p{1.6cm}|
>{\raggedright\arraybackslash}p{1.75cm}|
>{\raggedright\arraybackslash}p{2.45cm}|}
\hline
\textbf{Method} & \textbf{Avg. RSSI (dBm)} & \textbf{Avg. Time (s)} & \textbf{Avg. Moves} & \textbf{Success (\%)} & \textbf{RSSI Std. Dev. (dB)} \\
\hline
Exhaustive scan & $-48.95$ & 54.20 & 812 & 100 & 1.60 \\
\hline
Hill-climbing & $-52.73$ & 11.84 & 63 & 80 & 2.34 \\
\hline
RL (Q-learning) & $-49.55$ & $\mathbf{9.62}$ & $\mathbf{57}$ & $\mathbf{95}$ & $\mathbf{1.28}$ \\
\hline
\end{tabular}
\end{adjustbox}
\end{table}

As we can see from the total data obtained, reinforcement learning (RL) was much quicker compared to exhaustive searching while being able to produce similar RSSI and outperforming hill-climbing by other parameters. The exhaustively searching method had the highest RSSI value and 100\% success rate due to comprehensive exploring of the space. However, this high reliability was associated with a significant cost in time and moves. Hill-climbing proved to be faster than exhaustive searching but less successful and producing lower RSSI values.

\subsubsection{Scenario-wise Practical Comparison}

\vspace{0.5cm}
\renewcommand{\arraystretch}{1.25}
\setlength{\LTpre}{8pt}
\setlength{\LTpost}{8pt}
\begin{longtable}{|
>{\raggedright\arraybackslash}p{3.6cm}|
>{\raggedright\arraybackslash}p{3.5cm}|
>{\raggedright\arraybackslash}p{3cm}|
>{\raggedright\arraybackslash}p{3cm}|}
\caption{Scenario-wise Final RSSI Comparison}
\label{tab:scenario_rssi_comparison}\\
\hline
\textbf{Scenario} & \textbf{Exhaustive (dBm)} & \textbf{Hill-Climbing (dBm)} & \textbf{RL (dBm)} \\
\hline
\endfirsthead
\hline
\textbf{Scenario} & \textbf{Exhaustive (dBm)} & \textbf{Hill-Climbing (dBm)} & \textbf{RL (dBm)} \\
\hline
\endhead
Outdoor, 20~m & $-41.8 \pm 0.9$ & $-44.9 \pm 1.8$ & $-42.1 \pm 0.9$ \\
\hline
Outdoor, 30~m & $-47.1 \pm 1.1$ & $-50.7 \pm 2.4$ & $-47.6 \pm 1.1$ \\
\hline
Indoor, 10~m (LOS) & $-50.4 \pm 1.3$ & $-54.1 \pm 2.9$ & $-50.9 \pm 1.3$ \\
\hline
Indoor, 10~m (obstructed) & $-56.5 \pm 1.6$ & $-61.2 \pm 3.5$ & $-57.6 \pm 1.8$ \\
\hline
\end{longtable}

\vspace{0.5cm}
\renewcommand{\arraystretch}{1.25}
\begin{longtable}{|
>{\raggedright\arraybackslash}p{4.2cm}|
>{\raggedright\arraybackslash}p{3.0cm}|
>{\raggedright\arraybackslash}p{3.0cm}|
>{\raggedright\arraybackslash}p{3.0cm}|}
\caption{Scenario-wise Convergence Time and Success Rate}
\label{tab:scenario_time_success}\\
\hline
\textbf{Scenario} & \textbf{Exhaustive} & \textbf{Hill-Climbing} & \textbf{RL} \\
\hline
\endfirsthead
\hline
\textbf{Scenario} & \textbf{Exhaustive} & \textbf{Hill-Climbing} & \textbf{RL} \\
\hline
\endhead
Outdoor, 20~m & $49.2 \pm 1.5$ s, 100\% & $9.1 \pm 1.9$ s, 90\% & $\mathbf{7.8 \pm 1.2}$ s, 100\% \\
\hline
Outdoor, 30~m & $53.8 \pm 1.8$ s, 100\% & $10.8 \pm 2.2$ s, 80\% & $\mathbf{9.0 \pm 1.3}$ s, 100\% \\
\hline
Indoor, 10~m (LOS) & $55.1 \pm 2.0$ s, 100\% & $12.6 \pm 2.7$ s, 80\% & $\mathbf{10.1 \pm 1.6}$ s, 90\% \\
\hline
Indoor, 10~m (obstructed) & $58.7 \pm 2.3$ s, 100\% & $14.9 \pm 3.0$ s, 70\% & $\mathbf{11.6 \pm 1.8}$ s, 90\% \\
\hline
\end{longtable}

\vspace{0.5cm}
\begin{figure}[H]
\centering
\begin{tikzpicture}
\begin{axis}[
    ybar,
    bar width=7pt,
    width=\textwidth,
    height=7.2cm,
    ymin=-65,
    ymax=-35,
    ylabel={Final RSSI (dBm)},
    symbolic x coords={Out-20,Out-30,In-10L,In-10O},
    xtick=data,
    enlarge x limits=0.18,
    legend style={at={(0.5,1.05)}, anchor=south, legend columns=3},
    grid=major
]
\addplot[fill=gray!45] coordinates {(Out-20,-41.8) (Out-30,-47.1) (In-10L,-50.4) (In-10O,-56.5)};
\addplot[fill=orange!60] coordinates {(Out-20,-44.9) (Out-30,-50.7) (In-10L,-54.1) (In-10O,-61.2)};
\addplot[fill=blue!50] coordinates {(Out-20,-42.1) (Out-30,-47.6) (In-10L,-50.9) (In-10O,-57.6)};
\legend{Exhaustive,Hill-Climbing,RL}
\end{axis}
\end{tikzpicture}
\vspace{0.15cm}
\caption{Scenario-wise final RSSI comparison in practical trials}
\label{fig:practical_rssi_chart}
\end{figure}

\begin{figure}[H]
\centering
\begin{tikzpicture}
\begin{axis}[
    ybar,
    bar width=7pt,
    width=\textwidth,
    height=7.2cm,
    ymin=0,
    ylabel={Convergence Time (s)},
    symbolic x coords={Out-20,Out-30,In-10L,In-10O},
    xtick=data,
    enlarge x limits=0.18,
    legend style={at={(0.5,1.05)}, anchor=south, legend columns=3},
    grid=major
]
\addplot[fill=gray!45] coordinates {(Out-20,49.2) (Out-30,53.8) (In-10L,55.1) (In-10O,58.7)};
\addplot[fill=orange!60] coordinates {(Out-20,9.1) (Out-30,10.8) (In-10L,12.6) (In-10O,14.9)};
\addplot[fill=blue!50] coordinates {(Out-20,7.8) (Out-30,9.0) (In-10L,10.1) (In-10O,11.6)};
\legend{Exhaustive,Hill-Climbing,RL}
\end{axis}
\end{tikzpicture}
\vspace{0.15cm}
\caption{Scenario-wise convergence time comparison in practical trials}
\label{fig:practical_time_chart}
\end{figure}

Figures~\ref{fig:practical_rssi_chart} and~\ref{fig:practical_time_chart} show a
consistent practical trend across all four scenarios. Exhaustive scan achieved the
best final RSSI in each case, as expected from a full search of the sampled
orientation space. RL converged much faster than exhaustive scan while reaching a
final RSSI close to the exhaustive reference and still outperforming hill-climbing in
every scenario. The performance gap widened in the obstructed indoor 10~m case,
where the wireless channel was least smooth and hill-climbing was most vulnerable to
local maxima.

\vspace{0.5cm}
\begin{table}[H]
\centering
\fontsize{12}{14}\selectfont
\caption{Improvement of RL Relative to Baseline Methods}
\label{tab:rl_relative_improvement}
\renewcommand{\arraystretch}{1.25}
\begin{adjustbox}{max width=\textwidth}
\begin{tabular}{|
>{\raggedright\arraybackslash}p{3.8cm}|
>{\raggedright\arraybackslash}p{2.2cm}|
>{\raggedright\arraybackslash}p{2.2cm}|
>{\raggedright\arraybackslash}p{2.4cm}|
>{\raggedright\arraybackslash}p{2.4cm}|}
\hline
\textbf{Scenario} & \textbf{RSSI Gain vs Exh.} & \textbf{RSSI Gain vs Hill} & \textbf{Time Reduction vs Exh.} & \textbf{Time Reduction vs Hill} \\
\hline
Outdoor, 20~m & $-0.3$~dB & 2.8~dB & 84.1\% & 14.3\% \\
\hline
Outdoor, 30~m & $-0.5$~dB & 3.1~dB & 83.3\% & 16.7\% \\
\hline
Indoor, 10~m (LOS) & $-0.5$~dB & 3.2~dB & 81.7\% & 19.8\% \\
\hline
Indoor, 10~m (obstructed) & $-1.1$~dB & 3.6~dB & 80.2\% & 22.1\% \\
\hline
\end{tabular}
\end{adjustbox}
\end{table}

According to the table showing the improvement achieved, the RL algorithm never outperforms the exhaustive algorithm in terms of the final RSSI since it depicts the realistic result of an exhaustive search. However, the practical gain that can be achieved through this approach is the considerable reduction in the time it takes to converge, while at the same time only making a minor sacrifice concerning RSSI compared to exhaustive search, as well as an obvious superiority over hill climbing.

\subsection{Validation Through Theory, Simulation, and Practical Output}

The measured outputs agree well with the expected qualitative behavior of the three
algorithms. Table~\ref{tab:theory_sim_practical_validation} summarizes this validation.
The experimental results match very well the expected qualitative behaviors of the three algorithms.

\vspace{0.5cm}
\renewcommand{\arraystretch}{1.25}
\begin{longtable}{|
>{\raggedright\arraybackslash}p{2.35cm}|
>{\raggedright\arraybackslash}p{1.35cm}|
>{\raggedright\arraybackslash}p{2.5cm}|
>{\raggedright\arraybackslash}p{1.65cm}|
>{\raggedright\arraybackslash}p{1.35cm}|
>{\raggedright\arraybackslash}p{1.75cm}|
>{\raggedright\arraybackslash}p{1.65cm}|}
\caption{Simulation-to-Practical Deviation by Method}
\label{tab:simulation_practical_deviation}\\
\hline
\textbf{Method} & \textbf{Sim. RSSI} & \textbf{Prac. Avg. RSSI} & \textbf{$\Delta$ RSSI} & \textbf{Sim. Time} & \textbf{Prac. Avg. Time} & \textbf{$\Delta$ Time} \\
\hline
\endfirsthead
\hline
\textbf{Method} & \textbf{Sim. RSSI} & \textbf{Prac. Avg. RSSI} & \textbf{$\Delta$ RSSI (dB)} & \textbf{Sim. Time} & \textbf{Prac. Avg. Time} & \textbf{$\Delta$ Time} \\
\hline
\endhead
Exhaustive scan & $-30.14$ & $-48.95$ & $-18.75$ & 67.68 s & 54.20 s & $-19.9\%$ \\
\hline
Hill-climbing & $-32.13$ & $-52.73$ & $-22.60$ & 6.51 s & 11.84 s & $+81.9\%$ \\
\hline
RL (Q-learning) & $-30.82$ & $-49.55$ & $-19.73$ & 5.51 s & 9.62 s & $+74.6\%$ \\
\hline
\end{longtable}

In practice, the actual RSSI values have significantly decreased relative to the simulation in all three cases, showing that there is indeed a distinction between simulations and real-world scenarios. For the actual experiment, there are some aspects causing poorer readings due to misaligned antennas and mismatched radiation patterns, RF losses in wires and connectors, and the slip ring itself, low resolution of RSSI per packet measured by the ESP32 microcontroller, and ambient noise. All these aspects contribute to poorer RSSIs than the simple simulator, but the difference between the two is more obvious when hill climbing is used because of the large number of steps required. Nevertheless, despite reinforcement learning having an edge over hill climbing regarding the last RSSI and consistency, the actual experiments show that its performance is still inferior to the exhaustive reference.

\vspace{0.5cm}
\renewcommand{\arraystretch}{1.25}
\begin{longtable}{|
>{\raggedright\arraybackslash}p{3.6cm}|
>{\raggedright\arraybackslash}p{4.0cm}|
>{\raggedright\arraybackslash}p{2.95cm}|
>{\raggedright\arraybackslash}p{3.0cm}|}
\caption{Validation of Observed Behavior Against Expected Theory}
\label{tab:theory_sim_practical_validation}\\
\hline
\textbf{Expected Behavior} & \textbf{Theory} & \textbf{Simulation Output} & \textbf{Practical Output} \\
\hline
\endfirsthead
\hline
\textbf{Expected Behavior} & \textbf{Theory} & \textbf{Simulation Output} & \textbf{Practical Output} \\
\hline
\endhead
Exhaustive scan should be reliable but slow & Full search should locate the best sampled orientation at the cost of many measurements & Highest convergence time and largest step count & 100\% success in all scenarios, but by far the slowest method \\
\hline
Hill-climbing should be fast but locally sensitive & Greedy local search can terminate at suboptimal peaks & Fast convergence but low-to-moderate robustness & Short practical convergence time but reduced success rate and weaker final RSSI indoors \\
\hline
RL should balance speed and robustness & Learned value estimates should guide the search beyond immediate local RSSI gradients & Fastest simulated convergence with competitive final RSSI & Near-exhaustive average RSSI, faster convergence than exhaustive scan, and strong repeatability \\
\hline
Indoor channels should be harder than outdoor channels & Multipath and obstruction distort the RSSI surface & Not fully captured in idealized form & Lower final RSSI and longer convergence time for all methods, especially hill-climbing \\
\hline
\end{longtable}

\subsection{Limitations of RL}

Though the RL controller performed well in all test cases considered, there are certain limitations to note about the chosen Q-learning approach.

\textbf{Dependence on discretization:}
The implemented controller works with discrete values for pan, tilt, and RSSI trend states. Such discretization ensures the Q-table remains small enough to fit into the memory of the ESP32-S3 board; however, it confines the learned policy to a set of preselected orientations, possibly overlooking better configurations lying between the chosen angles.

\textbf{Generalization limits:}
Since the controller is based on a fixed Q-table obtained offline, its generalization ability is relatively lower compared to other types of models. Excellent performance in both indoor and outdoor experiments does not guarantee similar results when using different antenna designs, search spaces, motion precision, or transmission channels.

\textbf{Sensitivity to the training environment:}
The learned policy includes the characteristics of the used simulator and offline training process. In case of increased multipath interference, additional attenuation, fluctuation of the received signal strength, or mechanical inaccuracies in the real scenario, the actual policy deviates from simulated instructions, leading to lower RSSI levels and increased runtime.

\subsection{Error Analysis}

The difference between expected and measured outputs can be traced to several practical
non-idealities in the hardware platform and test environment.

\vspace{0.5cm}
\renewcommand{\arraystretch}{1.25}
\begin{longtable}{|
>{\raggedright\arraybackslash}p{3cm}|
>{\raggedright\arraybackslash}p{3.2cm}|
>{\raggedright\arraybackslash}p{3cm}|
>{\raggedright\arraybackslash}p{4.45cm}|}
\caption{Sources of Error and Their Impact on Results}
\label{tab:error_analysis}\\
\hline
\textbf{Source of Error} & \textbf{Effect on Output} & \textbf{Observed Evidence} & \textbf{Mitigation Used} \\
\hline
\endfirsthead
\hline
\textbf{Source of Error} & \textbf{Effect on Output} & \textbf{Observed Evidence} & \textbf{Mitigation Used} \\
\hline
\endhead
Antenna mismatch and alignment tolerance & Shifts the practical radiation pattern away from the idealized simulated pattern and reduces achievable peak RSSI & Practical peak RSSI remained below the simulated reference even after convergence & Careful mechanical assembly, fixed mounting geometry, and repeated trials \\
\hline
RF losses in connectors, cables, and slip-ring paths & Introduces attenuation and small signal variation not present in the simplified simulator & Consistent negative RSSI offset between simulation and practical measurements across all methods & Short signal paths, local decoupling, and separation of power and signal routing \\
\hline
ESP32 RSSI reporting limitations & Packet-level RSSI has finite resolution and short-term fluctuation, which affects reward quality and state estimation & Trial-to-trial RSSI variation remained even at similar orientations & RSSI averaging, trend-based state encoding, and fixed sampling windows \\
\hline
Multipath propagation, obstruction, and environmental noise & Creates local RSSI peaks and fluctuating measurements that distort the search landscape & Larger variance and poorer hill-climbing performance in indoor and obstructed cases & RSSI averaging, repeated trials, comparison across indoor and outdoor scenarios \\
\hline
Mechanical backlash and finite step resolution & Causes the realized angle to differ slightly from the commanded angle & Practical RSSI lower than simulated ideal values and slower settling after movement & Limited search range, calibrated step factor, mechanical settling delay \\
\hline
Slip-ring contact variation and electrical noise & Perturbs supply quality and signal stability during rotation & Trial-to-trial RSSI fluctuation and small convergence variability & Local decoupling, signal separation, short averaging window before decision making \\
\hline
Servo positioning error and step quantization & Prevents exact arrival at the continuous optimum orientation & Residual difference between exhaustive reference and final aligned angle & Discrete state/action design matched to hardware resolution \\
\hline
Simulation model simplification & Underestimates real environmental complexity & All practical RSSI values are lower than simulated values & Practical validation used as final evidence, not simulation alone \\
\hline
\end{longtable}

Amongst these issues, the effect of environmental multipath, RF losses, limitations in RSSI measurement, and actuator imperfections had the largest impact on the difference from expected behavior. It can also be argued that this list represents the main reason why hill climbing performs poorly indoors. The RL controller was less affected by these problems due to the fact that it is not dependent purely on the local RSSI gradient, but these issues still negatively impacted its performance compared to simulations.


\subsection{Overall Analysis}

Overall, from the engineering perspective, the results presented in this chapter show that the developed solution is viable in simulation and in the real environment as well. Since the Q-table used by the RL algorithm is relatively small, it is possible to implement the controller on the ESP32-S3 and test its convergence in practice. In such conditions, the RL controller provided the most optimal balance between convergence speed, RSSI value, and repeatability out of all the alternatives considered.

Practical results obtained during this study are important since they prove the feasibility of the solution in actual conditions, even though the simulation-real difference and limitations inherent to the used tabular RL framework have to be taken into account when interpreting the results obtained.